\documentclass{article}
\usepackage{amsmath}

\title{Convergence Behavior of Constraint-Based Physics Engines}
\author{Hyunseung Ryu}

\begin{document}
\maketitle

\begin{abstract}
Position-based dynamics trades physical accuracy for stability, and the trade is
governed by a single knob: how many times the solver sweeps the constraint set.
We measure how that count interacts with stiffness and show that the effective
stiffness saturates well before the iteration budget does.
\end{abstract}

\section{Setup}

Let $\mathbf{x}_i$ denote particle positions and $w_i = 1/m_i$ the inverse masses.
A constraint $C(\mathbf{x})$ is projected by moving each particle along the
gradient direction, weighted by its inverse mass:

\begin{equation}
  \Delta \mathbf{x}_i = -\frac{w_i \, \nabla_{\mathbf{x}_i} C(\mathbf{x})}
    {\sum_j w_j \left\| \nabla_{\mathbf{x}_j} C(\mathbf{x}) \right\|^2} \, C(\mathbf{x}).
\end{equation}

The formulation follows \cite{muller2007}, which introduced the projection scheme
this work builds on.

\subsection{Stiffness and iteration count}

Applying a stiffness factor $k \in [0, 1]$ once per iteration gives an effective
stiffness after $n$ sweeps of

\begin{equation}
  k' = 1 - (1 - k)^{n}.
\end{equation}

The dependence on $n$ is the problem: the same $k$ produces a stiffer material
simply because the solver ran longer. The correction proposed by \cite{macklin2016}
removes the coupling.

\section{Results}

\begin{itemize}
  \item Small $k$ converges slowly but stays stable under large time steps.
  \item Large $k$ overshoots, and the overshoot grows with $n$.
  \item Past $n \approx 20$ the residual falls below the integration error.
\end{itemize}

\section{Conclusion}

Iteration count should not be a material parameter. Decoupling it, as in
\cite{macklin2016}, keeps the simulation reproducible across solver budgets.

\end{document}
